\documentclass[12pt]{letter}
\usepackage[left=1in,right=1in,top=1in,bottom=1in]{geometry}
\usepackage{lmodern}
\usepackage[T1]{fontenc}
\usepackage{setspace}
\onehalfspacing

\signature{Darcio Genicolo-Martins\\Paulo Furquim de Azevedo\\INSPER, S\~ao Paulo, Brazil}
\address{INSPER\\Rua Quat\'a, 300\\S\~ao Paulo, SP 04546-042\\Brazil}
\date{April 2026}

\begin{document}

\begin{letter}{The Editors\\Journal of Law and Economics\\University of Chicago}

\opening{Dear Editors,}

We submit ``Screening for Bid Rigging with Frequent Losers in Public Procurement'' for consideration at the \textit{Journal of Law and Economics}.

The paper introduces a participation-based screen for bid-rigging cartels. The key observation is simple: a firm that bids on dozens of procurement tenders over a decade without winning a single one is unlikely to be competing. We call these firms \textit{frequent losers} (FL) and show that they mark environments with systematically higher prices, deep proximity to known cartels, and bid patterns consistent with coordinated cover bidding.

The screen requires only win/loss records---no bid microdata, no enforcement priors, no supervised training. Applied to 4.5 million tender-items from Brazilian procurement, it achieves AUC~$= 0.94$ against CADE (Brazil's competition authority) cartel convictions. FL-flagged tenders show 3.6--7.7\% higher prices, and five independent diagnostics produce a joint pattern coherent with coordinated cover bidding.

We believe the paper is well suited for JLE for three reasons:

\begin{enumerate}
\item \textbf{Law enforcement relevance.} The screen is designed as a triage tool for competition authorities. We validate it against actual cartel convictions and propose a three-stage enforcement workflow (screen, triage, investigate). The oversight-heterogeneity result---a 12.5$\times$ gradient in the FL price gap between the smallest and largest purchasing units---speaks directly to the allocation of enforcement resources.

\item \textbf{Minimal data requirements.} The screen operates where bid-level forensics are infeasible, a common constraint in developing economies with weak procurement oversight. This extends the reach of cartel detection to settings where existing tools cannot be deployed.

\item \textbf{Institutional design implications.} The model predicts, and the data confirm, that the minimum-bidder rule in sealed-bid procurement (convite) creates a direct incentive for cover bidding. This has implications for procurement regulation, including Brazil's recent Lei 14.133/2021 reform.
\end{enumerate}

The manuscript has not been submitted elsewhere. A companion paper developing the full structural model and policy counterfactuals is available as a working paper. All data and code will be made available upon acceptance in compliance with JLE's replication policy.

\closing{Sincerely,}

\end{letter}
\end{document}
